% ===========================================================================
\section{Discussion}
\label{sec:discussion}
% ===========================================================================

\subsection{What T1 Buys, and What It Honestly Does Not}

The captured run demonstrates the T1 claim precisely: after the device
signs, every byte of evidence is bound to the ledger, and single-sample
tampering is detected by two independent mechanisms.  It equally
demonstrates the boundary: had the bridge host synthesized the entire
115-sample trace before signing, nothing downstream could tell.  We regard
stating this boundary as a contribution in itself.  A certifying authority
needs to know that a T1 record means ``this evidence has not changed since
the device signed it,'' not ``this print happened as described.''  For the
dominant practical threats---retroactive editing when a part later fails,
reuse of good evidence, quiet design substitution---tamper-evidence after
signing is exactly the needed property, deployable today on unmodified
consumer hardware.

\subsection{Path to T2 and T3}

T2 requires no new cryptography: the device identity and raw public key
already exist in the bundle, and the chaincode gains a registry keyed by
device ID with enrollment, revocation, and endorsement-time Ed25519
verification of the bundle.  Its value is organizational: enrollment
becomes a ledger event subject to endorsement policy, so a stolen key is
bounded by its revocation window and a bureau cannot quietly add signing
devices.  T3 addresses fabrication at source---key in a TEE or secure
element, measured boot of the recorder, and, on industrial machines,
signing at the sensor or controller level.  Threshold attestation across
independent observers (e.g., a bureau device plus a customer-supplied
witness camera or power monitor) raises collusion cost without trusted
hardware on the printer.  Even T3, honestly, does not defeat
physical-world spoofing (a detached thermistor in a heated block):
attestation raises the cost and forensic visibility of fraud; it does not
make physics self-authenticating.

\subsection{Assessment-Policy Calibration}

The quality score of 0.2 for the cancelled run, whose steady-state
tracking was tight ($\pm 0.5\,^{\circ}$C at setpoint), shows the current
policy is miscalibrated for short runs: idle and ramp samples, during which
targets are legitimately unmet, dominate the in-tolerance fraction.  The
five-run data makes the effect precise: the two idle-started cubes score
0.78--0.71 despite excellent steady-state tracking, whereas
\icode{cube\_repeat}, captured after warm-up, scores 0.94---the
$\approx$0.16 spread is almost entirely the heat-up artifact, not a
difference in print fidelity.  The fix is phase-aware scoring (weighting
only printing-phase samples), but we report the uncalibrated numbers
because the methodological point stands regardless: the policy is
transparent and deterministic, encoded in the signed bundle itself, so any
party can recompute and contest a rating from the evidence---unlike
subjective ratings, which are unfalsifiable.

\subsection{Limitations}

(i)~The recorder polls OctoPrint's REST API rather than hooking its event
bus, adding one software layer inside the T1 trust domain.  (ii)~A 5\,s
poll can miss sub-interval transients; the thermal-anomaly run of
Section~\ref{sec:eval_anomaly} captured a 60\,s excursion comfortably, but
faster transients would require event-bus hooking or a shorter interval.
(iii)~The evaluation is single-printer, five-job, three-geometry: the
remaining gap is breadth across \emph{machines and sites}, not across jobs
or part shapes.  One of the five runs (\icode{cube\_repeat}) had its
recorder attach after warm-up, so we use it only as a completed-print data
point and reserve the controlled quality comparison for the two
idle-started cube runs---a reminder that capture must begin before heating
for heat-up samples to be scored consistently.  (iv)~Outcome
classification cannot separate cancellation from failure; the event-hook
plugin can.  (v)~Bundle growth is linear in duration; summarization
policies that retain raw extremes while downsampling steady state are
future work.  (vi)~A read-only ledger query
(\icode{GetAssetHistory}) currently fails an output-schema check
(\icode{linkedRatingID}) on history events; this is a query-path defect
only---all write-path commitments and the reputation updates of
Section~\ref{sec:eval_reputation} succeed and verify---but it should be
fixed before the trust-report view is relied upon operationally.
